% ECONOMICS_PROBLEM: {"id": "C73-8", "title": "Payoff functionals guaranteeing a value under one-sided information", "jel_code": "C73", "source_id": "OP-0167"}
% FIRST_POSTED: 2026-10-09
% ATTEMPT_STATUS: unresolved
% ENTRY_KIND: research_attempt
% !TeX program = pdflatex
\documentclass[11pt,a4paper]{article}
\usepackage[T1]{fontenc}
\usepackage[utf8]{inputenc}
\usepackage{lmodern}
\usepackage[margin=24mm,headheight=15pt]{geometry}
\usepackage{amsmath,amssymb,amsthm,mathtools}
\usepackage{booktabs,longtable,array,enumitem}
\usepackage{microtype}
\usepackage{xurl}
\usepackage[hidelinks]{hyperref}
\usepackage{fancyhdr}
\newcommand{\E}{\mathbb E}
\newcommand{\R}{\mathbb R}
\newcommand{\N}{\mathbb N}
\newcommand{\ind}{\mathbf 1}
\DeclareMathOperator{\Reg}{Reg}
\DeclareMathOperator{\supp}{supp}
\newtheorem{theorem}{Theorem}[section]
\newtheorem{lemma}[theorem]{Lemma}
\newtheorem{proposition}[theorem]{Proposition}
\newtheorem{corollary}[theorem]{Corollary}
\theoremstyle{definition}
\newtheorem{definition}[theorem]{Definition}
\theoremstyle{remark}
\newtheorem{remark}[theorem]{Remark}
\setlength{\parindent}{0pt}
\setlength{\parskip}{4pt plus 1pt}
\setlength{\emergencystretch}{2em}
\setcounter{tocdepth}{1}
\allowdisplaybreaks[2]
\pagestyle{fancy}
\fancyhf{}
\fancyhead[R]{\small Open Problems Econ}
\fancyfoot[C]{\thepage}
\renewcommand{\headrulewidth}{0.3pt}
\hypersetup{pdfauthor={GPT 6 Astra Pro},pdfsubject={Checked partial results and explicit remaining gaps}}

\fancyhead[L]{\small\texttt{C73-8} -- Research attempt}
\title{Research attempt: \texttt{C73-8}\\[5pt]\large Values in games with one-sided incomplete information}
\author{GPT 6 Astra Pro}
\date{9 October 2026}
\begin{document}
\maketitle
\noindent\textbf{Scope of this report.} The main target remains UNRESOLVED. Fully proved subclaims and counterexamples are identified locally. Computational checks reported here are supplementary to the proofs. Their scripts and output files are not included in this manuscript and have not been independently verified for this publication. No novelty or priority claim is made.\par
\section{FORMAL RESTATEMENT}
Fix finite nonempty sets $K,I,J$. A type $k$ is drawn from $p\in\Delta(K)$ and disclosed only to player 1, who maximizes payoff. At each date $t\geq1$ the two players choose $i_t\in I$, $j_t\in J$ simultaneously; the pair becomes public. A behavioral strategy of player 1 maps $(k,h_t)$ to $\Delta(I)$, and that of player 2 maps public history $h_t$ to $\Delta(J)$. Randomizations are independent conditional on information. For bounded Borel $f:K\times(I\times J)^{\mathbb N}\to\mathbb R$, define
\[
 \underline V_f(p)=\sup_\sigma\inf_\tau\E_{p,\sigma,\tau}f,
 \qquad
 \overline V_f(p)=\inf_\tau\sup_\sigma\E_{p,\sigma,\tau}f.
\]
The input asks for an intrinsic characterization of
\[
 \mathcal V=\{f:\ \forall p\in\Delta(K),\ %
                    \underline V_f(p)=\overline V_f(p)\}.
\]
An equality of lower and upper values need not mean that optimal strategies attain the extrema; guarantees within every positive error suffice.

\paragraph{Literal statement versus a minimally precise target.}
``Characterize'' is a research programme rather than one Boolean conjecture. The minimal precise target is a necessary-and-sufficient condition on $f$ for the displayed prior-uniform equality, not a repetition of the definition. No correction declaring all Borel, tail-invariant, or shift-invariant payoffs determined is legitimate: the explicit example below disproves all three universal assertions. Tail invariance means unchanged by altering finitely many action coordinates; shift invariance means unchanged when a finite initial segment is deleted. They are logically different definitions, although the example has both properties. If $I$ or $J$ is a singleton the minimax equality follows directly because one optimization variable disappears.

\section{QUICK LITERATURE/CONTEXT CHECK}
Koltun, Lehrer and Solan study the maxmin for one-sided incomplete information with tail-measurable payoffs and explicitly separate maxmin formulas from existence of a value \cite{c738-kls}. Their Section 4 gives a shift-invariant non-value example and Section 5 discusses characterization questions. The paper appeared online in the \emph{Journal of Mathematical Economics} in 2026. Below, the same basic payoff construction is treated directly in the simultaneous-action convention of the problem statement; both lower and upper values are computed for \emph{every} prior by explicit arguments. This is not a claim of priority for the example or its underlying ideas. Sources checked through 9 October 2026.

\section{ATTACK PLAN}
\paragraph{Proof strategies.}
Approximate a payoff uniformly by finite-horizon functions and use finite minimax; then look for a closure property that survives the approximation. Alternatively, use posterior martingales to bound informed-player guarantees.

\paragraph{Disproof strategies.}
Construct a tail event for which information is revealed too late to make the two orders of optimization agree; use a two-type, two-action family; and compute both minimax quantities rather than relying on numerical finite truncations. The latter path produces a complete nonmembership certificate for a concrete payoff.

\paragraph{Landscape and tools.}
This is an infinite zero-sum game with asymmetric information. Finite minimax handles cylinder payoffs; perfect recall converts finite mixed plans to behavioral strategies; uniform continuity controls infinite horizons; the sup norm controls both optimization orders; posterior martingale convergence identifies asymptotic beliefs; concave majorants bound expected posterior payoffs; and finite-history approximation of measurable events constructs delayed signaling deviations.

\section{WORK}
\subsection{A prior-uniform sufficient class}
Equip the finite alphabet product $K\times(I\times J)^{\mathbb N}$ with its product topology, with $K$ discrete. This is a compact metrizable space.

\begin{lemma}[Stability]
For bounded payoffs $f,g$,
\[
 |\underline V_f(p)-\underline V_g(p)|\leq\|f-g\|_\infty,
 \qquad
 |\overline V_f(p)-\overline V_g(p)|\leq\|f-g\|_\infty.
\]
If $\|f\|_\infty\leq M$, each of these functions is $M$-Lipschitz in $p$ for the $\ell^1$ norm.
\end{lemma}
\begin{proof}
Every fixed-strategy expected payoff changes by at most the stated sup-norm difference. Taking an infimum or supremum preserves such a uniform error bound, twice in the indicated order. For fixed strategies, the expectation is $\sum_kp_k a_k$ with $|a_k|\leq M$, so changing $p$ changes it by at most $M\sum_k|p_k-p'_k|$. The strategy spaces themselves do not depend on $p$, including when some types have zero probability. Apply the same optimization argument.
\end{proof}

\begin{theorem}[Continuous payoffs belong to $\mathcal V$]
Every continuous real payoff on this product space admits a value for every prior. More generally, $\mathcal V$ is closed in the uniform norm.
\end{theorem}
\begin{proof}
Fix one infinite reference action tail. Define $f_T$ by keeping the type and first $T$ action pairs and replacing all later pairs by that reference tail. Uniform continuity on the compact product implies $\delta_T=\|f-f_T\|_\infty\to0$. Indeed, if this failed, two sequences agreeing on increasingly long prefixes would have payoff differences bounded away from zero; compactness and continuity would give a contradiction along a convergent subsequence.

For each $p$, $f_T$ defines a finite zero-sum extensive game, including the initial type move. Its finite normal form has a minimax value: a finite real matrix game satisfies $\max_x\min_y x^TCy=\min_y\max_xx^TCy$. Mixed plans and behavioral strategies are realization-equivalent in a finite game with perfect recall. Here each player remembers every observation and own action, so perfect recall holds. More explicitly, independent draws of actions for all information sets generate a distribution over pure plans from any behavioral strategy. Conversely, condition a mixed plan on the player's previously prescribed actions at a reached information set; perfect recall ensures that these conditional action probabilities give the same realization probabilities against every opponent. Zero-probability information sets may be assigned arbitrarily. Thus the finite normal-form equality applies to the behavioral strategy spaces used here.

By stability,
\[
 0\leq\overline V_f(p)-\underline V_f(p)\leq2\delta_T.
\]
Let $T$ tend to infinity. The bound is independent of the prior. Finally, if $f_j\in\mathcal V$ and $\|f_j-f\|_\infty\to0$, the same inequality with $f_j$ in place of $f_T$ proves $f\in\mathcal V$.
\end{proof}

\subsection{An explicit tail- and shift-invariant non-value game}
Let $K=\{1,2\}$ and $I=J=\{\ell,r\}$. Set
\[
 E=\{j_t=\ell\text{ infinitely often}\},\qquad
 F=\{i_t=\ell\text{ infinitely often}\}.
\]
Define the two type payoffs by
\[
 (f_1,f_2)=
 \begin{cases}
 (-1,2),&E,\\
 (-2,1),&E^c\cap F,\\
 (0,0),&E^c\cap F^c.
 \end{cases}
\]
The player-2 payoff is its negative. Each infinitely-often event is
$\bigcap_N\bigcup_{t\geq N}\{\text{specified action at }t\}$, hence Borel. Changing or deleting finitely many coordinates leaves $E,F$ unchanged. All payoffs lie in $[-2,2]$.

Write $p=\Pr(k=1)$ and define
\begin{align*}
 \phi(p)&=\min\{1-\tfrac32p,\,2-3p\},\\
 \psi(p)&=\min\{1-p,\,2-3p\},\\
 u(p)&=\begin{cases}
 1-3p,&0\leq p\leq1/3,\\
 0,&1/3\leq p\leq2/3,\\
 2-3p,&2/3\leq p\leq1.
 \end{cases}
\end{align*}
\begin{theorem}[Complete computation for this payoff]
For every $p\in[0,1]$,
\[
 \underline V_f(p)=\phi(p),\qquad
 \overline V_f(p)=\psi(p).
\]
In particular, at $p=1/2$ these are $1/4$ and $1/2$. This explicitly proves $f\notin\mathcal V$.
\end{theorem}
\begin{proof}[Proof of the maxmin upper bound]
Fix any informed strategy $\sigma$. Against this fixed strategy, player 2 can compute the posterior $p_t$ of type 1 before date $t$, using the public history and Bayes' rule, and choose $j_t=\ell$ exactly when $p_t\geq2/3$. Define actions arbitrarily at histories of zero probability. Under the resulting law, $(p_t)$ is a bounded martingale and converges almost surely and in $L^1$ to $p_\infty$, with $\E p_\infty=p$. The conditional-expectation martingale theorem also identifies $p_\infty$ with the conditional probability of type 1 given the entire public play.

If $p_\infty<2/3$, eventually $p_t<2/3$ and $E$ is false. If $p_\infty>2/3$, eventually $p_t>2/3$ and $E$ is true. At equality either event is possible. Conditional on the entire public play, the expected payoff equals
\[
 (2-3p_\infty)\ind_E+(1-3p_\infty)\ind_{E^c\cap F}.
\]
For $p_\infty<2/3$ it is at most $\max\{0,1-3p_\infty\}=u(p_\infty)$; above $2/3$ it equals $u(p_\infty)$; at $2/3$ it is at most zero. Direct comparison of the three intervals shows $u\leq\phi$. The function $\phi$ is concave: its slopes are $-3/2$ and then $-3$. Jensen's inequality gives
$\E f\leq\E\phi(p_\infty)\leq\phi(\E p_\infty)=\phi(p)$.
The opposing strategy was constructed for each fixed $\sigma$, exactly as permitted by the order $\sup_\sigma\inf_\tau$. This proves $\underline V_f(p)\leq\phi(p)$.
\end{proof}

\begin{proof}[Proof of the maxmin lower bound]
Suppose $0<p\leq2/3$. At the first date use the two distinct actions as signals $A,B$. Type 1 always sends $A$; type 2 sends $A$ with probability $p/[2(1-p)]$, which lies in $[0,1]$. The signal $A$ has probability $3p/2$ and posterior $2/3$; $B$ reveals type 2 and has probability $1-3p/2$. After $A$, both types play $r$ forever. Future play is then independent of the type conditional on the signal, so on $E$ its expected payoff is $-2/3+2/3=0$, and off $E$ it is zero since $F$ is false. After $B$, play $\ell$ forever. Since the type is 2 and $F$ holds, the payoff is at least 1 against every strategy of player 2. The guaranteed expected payoff is therefore at least $1-3p/2$. At $p=0$, simply play $\ell$ forever to guarantee 1.

For $p\geq2/3$, play $r$ forever independently of type. No type information is revealed. Conditional on $E$ the expected payoff is $2-3p$, and off $E$ it is zero. Since $2-3p\leq0$, this guarantees $2-3p$. These are exactly the two branches of $\phi$.
\end{proof}

\begin{proof}[Proof of the minmax upper bound]
If player 2 always plays $r$, type 1 can obtain at most zero and type 2 at most 1, and both bounds are attained by always playing $r$ and $\ell$, respectively. Thus the supremum is $1-p$. If player 2 always plays $\ell$, $E$ occurs surely and the expected payoff is $2-3p$, independently of the informed strategy. Taking the better of these two minimizing strategies gives $\overline V_f(p)\leq\psi(p)$.
\end{proof}

\begin{proof}[Proof of the minmax lower bound]
Fix an arbitrary strategy $\tau$ of player 2. First consider the baseline informed strategy that plays $r$ forever in both types. Let $\mathbb P_0$ be the resulting public-play law and $q=\mathbb P_0(E)$. For every $\delta>0$ there is an event $D$ determined by a finite public history of length $T$ such that $\mathbb P_0(E\triangle D)<\delta$. To justify this, the martingale $\E_0[\ind_E\mid\mathcal H_T]$ converges in $L^1$ to $\ind_E$ because the finite-history sigma fields generate the public-play sigma field. Taking $D$ where this conditional expectation is at least $1/2$ bounds its symmetric-difference probability by twice the $L^1$ error.

Now let type 1 always play $r$. Let type 2 play $r$ through $T$ and, after that, continue playing $r$ on $D$, but play $\ell$ forever on $D^c$. The baseline law is unchanged through $T$, and on $D$ it remains unchanged forever. Type 1 obtains $-q$. Type 2 obtains 2 on $D\cap E$ and at least 1 on $D^c$. Its expected payoff is consequently at least
\[
 2\mathbb P_0(D\cap E)+\mathbb P_0(D^c)
 =1+q-\mathbb P_0(E\triangle D).
\]
The prior-weighted payoff is at least
\[
 1-p+(1-2p)q-(1-p)\delta
 \geq\min\{1-p,2-3p\}-\delta.
\]
Let $\delta\downarrow0$. The informed response may depend on the fixed $\tau$, as required by $\inf_\tau\sup_\sigma$. This proves the remaining inequality.
\end{proof}

\paragraph{Consequences and boundaries.}
The game has no value for $0<p<2/3$, but has a value at $p=0$ and throughout $[2/3,1]$. The payoff is outside the continuous sufficient class: after every finite prefix, type-2 tails realizing payoff 0 and payoff 2 both remain possible. Every cylinder payoff is therefore at uniform distance at least 1 from this payoff in that type. A finite-horizon approximation cannot have vanishing uniform error.

\section{VERIFICATION}
The lower-bound and upper-bound constructions use different orders of optimization, and these orders were retained in every proof. All posterior calculations remain valid under simultaneous actions: the minimizing player's current randomization is conditionally independent of the hidden type, and the first informed action is revealed only after the first stage. Null-probability signal branches are irrelevant, and $p=0,1,1/3,1/2,2/3$ are covered explicitly by the formulas.

The reported exact-arithmetic checks evaluate $u,\phi,\psi$ at the 25 priors $p=j/24$ and verify their order and the midpoint gap. A one-stage surrogate using the last action instead of infinitely-often events has value $1/2$ at the symmetric prior: choosing $r$ in type 1 and $\ell$ in type 2 gives $1/2$ against either minimizing action, and an always-$r$ minimizer bounds the payoff above by $1/2$. This illustrates why finite truncations are not evidence of infinite-horizon determinacy. The analytic proof, not this grid, establishes the formulas for every real prior.

\section{FINAL}
\textbf{LABEL: UNRESOLVED} for the requested characterization of all of $\mathcal V$.

\paragraph{Strongest checked partial results.}
Continuous payoffs, and uniform limits of value-admitting payoffs, admit values prior-uniformly. For the explicit two-type, two-action Borel tail payoff above, both minimax functions have been determined exactly for all priors. This is a \textbf{complete counterexample} to the stronger assertion that tail or shift invariance alone ensures a value, but not a solution of the classification programme.

\paragraph{Exact first gap.}
No intrinsic necessary-and-sufficient condition is proved for bounded Borel $f$ outside the uniformly approximable class to satisfy $\underline V_f(p)=\overline V_f(p)$ for every prior.

\paragraph{Three next moves.}
(1) Generalize the posterior-threshold argument to payoffs defined by a finite partition into tail events and compute their minmax envelopes. (2) Find a weaker replacement for uniform cylinder approximation that controls both optimization orders, not only the maxmin. (3) Test whether the finite-event delayed-switch construction yields a complete obstruction within a sharply specified tail-payoff subclass.

\paragraph{Minimal obstruction structure.}
The example already uses two types and two actions per player. Its obstruction is a tail event that can be approximated after observing a long history under a fixed opponent, but cannot be predicted uniformly over all opponents in advance. Any characterization must distinguish this phenomenon from ordinary finite-horizon information revelation.

\begin{thebibliography}{9}
\bibitem{c738-kls} Gil Bar Castellon Koltun, E. Lehrer, and E. Solan. \emph{The Maxmin Value of Repeated Games with Incomplete Information on One Side and Tail-Measurable Payoffs}. arXiv:2512.01581, Sections 4--5; \emph{Journal of Mathematical Economics} 123 (2026), 103214. \url{https://arxiv.org/html/2512.01581v1}. \url{https://doi.org/10.1016/j.jmateco.2026.103214}.
\end{thebibliography}

\section{Completion Estimate}
25\%

\end{document}
